
Sampling atomistic configurations from the Boltzmann distribution with MC is trivial to understand and implement numerically. Begin with a starting structure $\bm{R}^{\mbox{\scriptsize{start}}}$, one take a step along the PES by randomly perturbing the atomic positions of $\bm{R}^{\mbox{\scriptsize{start}}}$ to generate a new structure $\bm{R}^{\mbox{\scriptsize{new}}}$, which will be either accepted (in this case, set $\bm{R}^{\mbox{\scriptsize{start}}}=\bm{R}^{\mbox{\scriptsize{new}}}$), or rejected (in which case we revert to $\bm{R}^{\mbox{\scriptsize{start}}}$) based on a pre-defined \emph{acceptance criteria}. The most common acceptance criterion is the Metropolis acceptance\cite{metropolis1953equation} [\cref{fig:sampling_methods}(a)], in which the probability of accepting $\bm{R}^{\mbox{\scriptsize{new}}}$ at temperature $T$ is given by:
%\begin{equation}
%    \label{eq:metropolis}
%    P(\bm{R^{\mbox{\scriptsize{new}}}})=\left\{\begin{array}{ll}
%        \exp\left[-\beta(E(\bm{R}^{\mbox{\scriptsize{new}}})-E(\bm{R}^{\mbox{\scriptsize{start}}}))\right] &  \mbox{if }  E(\bm{R}^{\mbox{\scriptsize{new}}})>E(\bm{R}^{\mbox{\scriptsize{old}}})\\
%        1 &\mbox{otherwise,} 
%    \end{array}
%    \right.
%\end{equation}
\begin{equation}
    \label{eq:metropolis}
    p(\bm{R}^{\mathrm{new}}|\bm{R})=\min\left[1,\frac{\pi(\bm{R}^{\mathrm{new}})}{\pi(\bm{R})}\right],
\end{equation}
which is written in the form of transition probability to be in line with the subsequent discussions on the probability theories for machine learning [\emph{c.f.} \cref{sec:prob_theory_fundamentals}], $\pi(\bm{R})$ is the probability of observing a structure $\bm{R}$ at a given temperature following \cref{eq:Boltzmann_distr},
and the new structure will be accepted if such a transition probability $p(\bm{R^{\mbox{\scriptsize{new}}}}|\bm{R})>\mbox{rand}(0,1)$. 
This procedure will be repeated for a large number of $N_{\max}$ steps to reach adequate sampling. It guarantees that the Markov chain is sampled from the equilibrium Boltzmann distribution, which corresponds to a canonical ensemble and satisfies the condition of detailed balance. To test whether MC sampling converges, one could run multiple MC samplings from different starting structures and compare the similarities of the final probability distributions acquired from these samplings, for instance, using the Kullback–Leibler divergence.\cite{kullback1951information} 

 The choice of the MC  moves plays an important role in enhancing the acceptance of new structures, thus increasing the efficiency of MC samplings for low energy structures. The basin-hopping\cite{wales1997global,wales1998global} [\cref{fig:sampling_methods}(d)] is  a widely adopted MC sampling approach in atomistic modelling, in which each MC move consists of a random perturbation of the atomic positions followed by a structural optimisation to the nearest local minimum on the PES. The local optimisation leads to a much lower $E(\bm{R}^{\mbox{\scriptsize{new}}})$, which correspondingly increases the acceptance probability. In this sense, instead of the full PES, basin hopping samples the reduced energy surface [straight lines in \cref{fig:sampling_methods}(d)]. Nevertheless, the design of the random perturbation in basin-hopping, which is typically system-specific, is still critical in ensuring we can achieve a high acceptance probability in MC sampling.  For this, the general guidance to design efficient MC moves is to follow the Bell-Evans-Polanyi principle\cite{bell1936theory,evans1938inertia} [\cref{fig:sampling_methods}(e)] which states that a lower (higher) $E(\bm{R}^{\mbox{\scriptsize{new}}})$ is typically separated from $E(\bm{R}^{\mbox{\scriptsize{start}}})$ with a lower (higher) energy barrier, hence taking a MC move that crosses the low energy barrier will tend to result in a higher acceptance probability. 

Here, we shall briefly review some examples of MC moves that have been applied in sampling the PES for chemical and material systems. The simplest move is the random perturbation on the Cartesian coordinates that translates the particle positions, which is effective for sampling the PES of systems bonded by a simple energy function, such as the Lennard-Jones clusters.\cite{wales1997global,wales1998global}  For flexible molecules, such as peptides, one could focus on perturbing the rotatable bonds independently while preserving the atomic connectivity.\cite{echenique2006definition} However, this will not work effectively for systems such as cyclic molecules in which the rotation of 4-body torsion angles is highly coupled with 3-body bond angle change, for which low-mode search\cite{kolossvary1996low} could be a better alternative. In this case, the low-energy vibrational eigenmodes  point to a set of orthogonal directions for atomic movements, along which the internal coordinates of a molecule could change cooperatively and optimally into a new conformation by avoiding steric clashes, following the Bell-Evans-Polanyi principle. However, deliberately biasing the MC move to cross high energy barriers could also be a desirable feature for exploring chemical reaction pathways. \cite{shang2013stochastic} 

In the context of solid state materials, the eigenvectors for the imaginary phonon modes point to the direction of phase transitions, which can be used to guide the exploration of the underlying PES.\cite{togo2013evolution,da2015phase}  In the case where the cell volume $V$ is allowed to change, the MC simulation will be performed under the $NPT$ ensemble with a fixed pressure $P$ (see later), in which the acceptance probability [\cref{eq:metropolis}] needs to be changed accordingly into:
\begin{equation*}
    P(\mbox{new})=\min\left\{(1,\exp\left[-\beta\left(\Delta E+P\Delta V-Nk_B T\ln \frac{V^{\mbox{\scriptsize{new}}}}{V^{\mbox{\scriptsize{start}}}} \right)\right]\right\}.
\end{equation*}
For multicomponent alloys, one could construct phase diagrams by performing samplings under the Variance-Constrained Semi-Grand Canonical (VCSGC) ensemble [\cref{fig:sampling_methods}(k)],\cite{aangqvist2019icet,sadigh2012scalable} which enables both the atomic positions and the chemical composition of the system to be `swapped' under certain constraints of the chemical potentials. In this case, the acceptance probability is formulated as:
\begin{equation*}
P(\mbox{new})=  \min\left\{1, \exp\left[-\beta \left( \Delta E + N \bar{\phi} \Delta c + N \bar{\kappa} (\Delta c)^2 \right) \right] \right\},
\end{equation*}
in which $\bar{\phi}$ and $\bar{\kappa}$ are dimensionless parameters controlling the average and variance of the composition, respectively. $\Delta c$ is the change in concentration resulting from the trial move. Finally, $N$ denotes the total number of particles in a system.
